\documentclass[a4paper,12pt]{article}

% Set margins
\usepackage[hmargin=2.5cm, vmargin=2.5cm]{geometry}

\frenchspacing

% Language packages
\usepackage[utf8]{inputenc}
\usepackage[T1]{fontenc}
% \usepackage[magyar]{babel}

% AMS
\usepackage{amssymb,amsmath}

% Graphic packages
\usepackage{graphicx}

% Colors
\usepackage{color}
\usepackage[usenames,dvipsnames]{xcolor}

% Enumeration
\usepackage{enumitem}

% Links
\usepackage{url}
\usepackage{hyperref}

\usepackage{fancyhdr}

\fancyhead[LO, LE]{Pi Mathematical, Physical, and Informatical Journal}
\fancyhead[RO, RE]{Vol. 1, Math., Scientific, 2024}

\usepackage{amsthm}
%\newtheorem{definition}{Definition}[section]
%\newtheorem{notation}{Notation}[section]
\newtheorem{theorem}{Theorem}[section]
\newtheorem{definition}[theorem]{Definition}
\newtheorem{notation}[theorem]{Notation}
\newtheorem{lemma}[theorem]{Lemma}
\newtheorem{proposition}[theorem]{Proposition}%[section]
\newtheorem{problem}[theorem]{Problem}%[section]
\newtheorem{corollary}[theorem]{Corollary}%[section]
\newtheorem{remark}[theorem]{Remark}%[section]
\newtheorem{example}[theorem]{Example}%[section]
\newtheorem{exercise}[theorem]{Exercise}%[section]
%\newtheorem{solution}[theorem]{Solution}%[section]
\newtheorem{conjecture}[theorem]{Conjecture}%[section]

\renewenvironment{proof}[1][]{\noindent\textit{Proof.} }{$\qed$}
\newenvironment{solution}[1][]{\noindent\textit{Solution.} }{$\qed$}

\begin{document}

\pagestyle{fancy}

\title{\sc Paper title \thanks{This research was supported by the \dots}}

% Please enter the authors' names, affiliations, and email addresses in the appropriate spaces
\author{
	\textbf{First Author} \\
	Affiliation 1 \\
	\href{mailto:first.author@email.com}{first.author@email.com}
\and
	\textbf{Second Author} \\
	Affiliation 2 \\
	\href{mailto:second.author@email.com}{second.author@email.com}
\and
	\textbf{Third Author} \\
	Affiliation 3 \\
	\href{mailto:third.author@email.com}{third.author@email.com}
}

% The date of the final version
\date{\textit{24 September, 2024}}

Pi Mathematical, Physical, and Informatical Journal

University of Miskolc

TODO: Use a banner here with logo!

\href{https://doi.org/10.1000/jmathpi-1}{DOI: 10.1000/jmathpi-????}

{\let\newpage\relax\maketitle}

\begin{abstract}
Text of the abstract.
\end{abstract}

{\def\thefootnote{}\footnotetext{Keywords: math, AMS Subject Classification: AB123, CD567}}

\section{Introduction}

Text of the Introduction.

\section{Title of the second section}

Please use the environments below:

\begin{definition}
Text of this definition.
\end{definition}

\begin{notation}
Text of this notation.
\end{notation}

\begin{theorem}\label{theorem:Key1}
Text of this theorem. Use the mathematical formulas in the text as $\frac{a}{b}$. 
\par 
\[
\text{Text of this math equation without counting}
\]
\begin{equation}\label{equation:Key2}
\text{Text of this math equation with counting, like } \dfrac{a}{b}
\end{equation}
Use the math formulas in the math environment as in equation \ref{equation:Key2}. 
\end{theorem}

\par By Theorem \ref{theorem:Key1}...

\begin{lemma}
Text of this lemma.
\end{lemma}

\begin{proposition}
Text of this proposition.
\end{proposition}

\begin{problem}
Text of this problem.
\end{problem}

\begin{corollary}
Text of this corollary.
\end{corollary}

\begin{remark}
Text of this remark.
\end{remark}

\begin{example}
Text of this example.
\end{example}

\begin{exercise}
Text of this exercise.
\end{exercise}

\begin{solution}
Text of this solution.
\end{solution}

\begin{conjecture}
Text of this conjecture.
\end{conjecture}

\begin{proof}
Text of this proof.
\end{proof}

\section{Title of the third section}

\begin{table}[!ht] % or [H] to turn off float
  \centering
  \begin{tabular}{|c|c|}
    \hline
    1 & 2 \\
    \hline
    3 & 4 \\
    \hline
  \end{tabular}
  \caption{Table caption.}\label{table:Key3}
\end{table}

By the Table \ref{table:Key3}\dots

\begin{figure}[!ht] % or [H] to turn off float
  \centering
  \includegraphics[width=3cm]{ME_Logo.png}
  \caption{Figure caption.}\label{figure:Key4}
\end{figure}

By the Figure \ref{figure:Key4}\dots

Use different kind of lists, for example
\begin{itemize}
\item unordered list,
\item ordered list,
\item hierarchical list.
\end{itemize}

\section{Title of the fourth section}%etc.

Refer to books \cite{erdos2013} and articles \cite{aczel1983}.

\paragraph{Acknowledgements.} Text of the acknowledgements.

\bibliographystyle{plaindin}
\bibliography{references}

\end{document}
